\chapter{Closed-Form $\alpha$ from the Golden Torus}
\label{ch:alpha_golden_torus}

\noindent The substrate-mechanism for $R \cdot r = 1/4$ closes at \textbf{Class B substrate-mechanism manifestation} via Axiom 4 self-saturation + Op14 Meissner-asymmetric coupling + a named substrate-mechanism identification step: the time-averaged phasor enclosed area at Axiom-4 saturation onset equals the Nyquist cell cross-section area, giving $\pi R r = \pi (d/2)^2 \Rightarrow R \cdot r = 1/4$ at $d = 1\,\ell_\text{node}$. The ropelength-minimality embedding-selection framing is sidestepped by this substrate-mechanism. Cross-particle universality (electron $(2,3)$ + proton $(2,5)$ + $\Delta$ baryon $(2,7)$) and cross-domain universality (electron-scale self-saturation + nuclear-scale Pd-D externally-driven saturation) both hold. Canonical anchors: \kbleaf{research/2026-05-31\_Q-EMBED-SEL-1\_step\_c\_result.md} (Class B closure); \kbleaf{research/2026-05-31\_Q-EMBED-SEL-1\_step\_c\_phase2\_cross\_particle\_result.md} (cross-particle); \kbleaf{research/2026-05-31\_Q-EMBED-SEL-1\_step\_c\_phase3\_cold\_fusion\_result.md} (cross-domain). \textbf{Class B caveat:} the named phasor-area-equals-Nyquist-cell-area identification is substrate-canonical INPUT (not Class 2 axiom-emergence from K4 $+$ Cosserat primitives alone). $\alpha$'s value $\alpha^{-1} = 4\pi^3 + \pi^2 + \pi$ is closed-form geometry whose \emph{scale} ($\sim 1/137$) is forced by the Compton-resonance trapping condition, but whose \emph{exact value} rests on ONE substrate-geometric identification per route ($R \cdot r = 1/4$ for the Golden-Torus route; $z_0 \leftarrow 1.187 \leftarrow p_c = 8\pi\alpha$ for the rigidity-percolation route) which the substrate does NOT independently select --- confirmed by four dynamic engine tests (cell-count, Cosserat-binding, dressed-eigenmode, chiral-dressing) that returned a negative selection record, the $S_{11}$-landscape-flatness finding, and the $z_0$ $\alpha$-circularity (the first-principles $z_0$-from-K4 route does not yield a parameter-free map to 137: exact $z_{\text{eff}} \to 6$ Maxwell--Calladine count, \texttt{alpha\_free\_map\_to\_137\_exists = False}).
It closes at \textbf{Class B substrate-mechanism manifestation via a named identification that is not independently derived} --- NOT ``derives'' and NOT a ``Zero-Parameter Closure.'' All three named $\alpha$-routes (dynamical selection, kinematic unit-bridge, rigidity-percolation) are closed-negative, so the $\alpha$-value is a \textbf{standing echo} --- the flip-condition stays live ($R \cdot r = 1/4$ forced without $\alpha$-circularity would flip it to a chord) and the route-space is not provably exhausted. A Class 2 lift would require deriving the phasor$\leftrightarrow$real-space area bijection from K4 $+$ Cosserat primitives alone; this is not achieved at present, so the closure stands at Class B. The chapter title (``Closed-Form $\alpha$ from the Golden Torus'') and \S Mathematical Closure framing read at this Class B substrate-mechanism manifestation level, with the named substrate-canonical identification as input.

\begin{objectivebox}
\begin{itemize}
    \item Close the AVE framework's logical loop in closed form: express the Fine Structure Constant $\alpha^{-1} = 4\pi^3 + \pi^2 + \pi$ as Golden-Torus geometry whose scale is forced by the Compton-resonance trapping condition. (The exact value rests on the one substrate-geometric identification $R \cdot r = 1/4$, which the substrate does not independently select --- Class B substrate-mechanism manifestation, not a first-principles ``derivation'' and not a zero-parameter closure; see the Class B caveat above.)
    \item Identify the electron natively as the $0_1$ unknot in real space, with a $(2,3)$ phase-space Clifford-torus winding pattern (2 d-axis windings, 3 q-axis windings), tensioned to dielectric ropelength on the discrete chiral lattice.
    \item Derive the Golden Torus geometry ($R = \varphi/2$, $r = (\varphi-1)/2$) from $S_{11}$ minimization (Universal Operator~6) plus Nyquist core-thickness constraint (Axiom~1).
    \item Evaluate the multipole decomposition on $\mathbb{T}^2 \subset S^3 \subset \mathbb{C}^2$ to obtain $\alpha^{-1}_{\text{ideal}} = 4\pi^3 + \pi^2 + \pi \approx 137.0363$ at $T \to 0$.
    \item Introduce the Vacuum Strain Coefficient $\delta_{\text{strain}}$ as the CMB-induced thermal correction bridging the cold asymptote to the measured CODATA value.
    \item State the falsifiable prediction: $\alpha$ runs with local thermal energy; the cold $137.0363$ is the absolute-zero mathematical asymptote.
\end{itemize}
\end{objectivebox}

\section{The Electron: $0_1$ Unknot with $(2,3)$ Phase-Space Winding}
\label{sec:electron_trefoil}

In standard particle physics, the electron is treated as a dimensionless point charge, leading to infinite self-energy paradoxes that require artificial renormalization. In AVE, the Electron ($e^-$) is identified natively as the ground-state topological defect of the chiral Laves K4 Cosserat crystal:

\begin{itemize}
    \item \textbf{Real-space topology:} the $0_1$ \emph{unknot} --- the simplest closed flux-tube loop, with no real-space crossings. The unknot's circumference is the reduced Compton wavelength $\hbar/(m_e c) = \ell_{node}$.
    \item \textbf{Phase-space topology:} a $(2, 3)$ Clifford-torus winding pattern in the bond-pair LC tank's $(V_{\text{inc}}, V_{\text{ref}})$ phasor space (2 windings on the d-axis, 3 windings on the q-axis). This is the smallest non-trivial coprime $(p, q)$ pair (lowest crossing number $c = 3$), forced by the uniqueness theorem on torus knots --- the electron is the lightest stable non-trivial topological soliton.
\end{itemize}

\textbf{The trefoil lives in phase space; the soliton lives in real space.} The unknot's real-space flux tube is tensioned by the vacuum's elastic Lattice Tension ($T_{max,g}$) to its absolute structural yield limit, while its phase-space $(V_{\text{inc}}, V_{\text{ref}})$ phasor trajectory closes as the $(2, 3)$ winding pattern on the Clifford torus. The Golden Torus radii $(R, r) = (\varphi/2, (\varphi-1)/2)$ derived below are the phase-space dimensions at which the phase-space winding closes consistent with the Nyquist, crossings, and screening constraints below. \emph{(Scope, 2026-06-24: ``closes'' here names the \textbf{static geometric embedding} of the $(2,3)$ winding on the Clifford torus --- the deformation-invariant Clifford-torus / Link texture the seed carries. It is \textbf{not} a dynamical/energetic ``closure'': the phase-space coupling-winding test found the $(2,3)$ does \textbf{not} re-emerge as a conserved dynamical time-orbit (NEGATIVE in both internal loci, \#415 + \#417); the dynamical orbit carries the carrier ratio, not the topological winding. Charge $= \mathrm{Link}(\partial\Omega, F) \in \mathbb{Z}$ stands as static topology. See \texttt{research/2026-06-24\_engine-reroute-epic-summary.md}.)}

\subsection{The Dielectric Ropelength Limit (The Golden Torus)}
\label{sec:golden_torus}

In a continuous mathematical space, a knotted tube can be shrunk infinitely small. However, because the substrate possesses a discrete minimum pitch (Axiom~1, Ch.~\ref{ch:fundamental_axioms}), a topological flux tube physically cannot be infinitely thin.

We define the phase-space winding pattern's geometry using the mathematical limits of \textbf{Ropelength}---the tightest a knotted phase-space trajectory can be pulled before its own minimum discrete thickness prevents further tightening. The immense elastic Lattice Tension ($T_{max,g}$) of the vacuum constantly seeks to minimize the stored inductive energy of the defect, tightening the $(2, 3)$ phase-space winding as compactly as physically possible. This tightening is violently halted by three rigid hardware bounding limits:

\begin{enumerate}
    \item \textbf{The Core Thickness ($d$):} The absolute minimum physical width of a propagating flux tube is exactly one fundamental lattice pitch. Normalized to the hardware grid, the fundamental diameter of the tube is rigidly locked at $d \equiv 1 \, \ell_{node}$.

    \item \textbf{The Self-Avoidance Constraint ($R - r = 1/2$):} As the knot pulls tight, the internal strands passing through the central hole of the torus compress against each other. To prevent the flux lines from attempting to occupy the exact same discrete node (which would trigger catastrophic dielectric rupture), the distance between their centerlines must be at least the tube diameter ($d=1$). For a torus knot, the closest geometric approach of the strands is $2(R-r)$. The physical packing limit structurally enforces $2(R-r) = 1 \implies R - r = 1/2$.

    \item \textbf{The Clifford-Torus Screening Limit ($R \cdot r = 1/4$):} The substrate-mechanism for the $R \cdot r = 1/4$ relation closes via Axiom~4 self-saturation $+$ Op14 Meissner-asymmetric coupling at the $(2,3)$ chirality $+$ a named substrate-mechanism identification step: the time-averaged phasor enclosed area at Axiom-4 self-saturation onset equals the Nyquist cell cross-section area (canonical at \kbleaf{research/2026-05-31\_Q-EMBED-SEL-1\_step\_c\_result.md} \S2.3). This gives $\pi R r = \pi(d/2)^2 \Rightarrow R \cdot r = 1/4$ at $d = 1\,\ell_\text{node}$. The $\Lambda_{surf} = \pi^2$ value follows from this substrate-native Axiom-4 self-saturation chain (see \S\ref{sec:pi_squared_rigor} below).
\end{enumerate}

\subsection{Substrate-Mechanism Derivation of $\Lambda_{\text{surf}} = \pi^2$ (Axiom-4 Self-Saturation)}
\label{sec:pi_squared_rigor}

The normalization $\Lambda_{\text{surf}} = \pi^2$ is substrate-derived from the Axiom-4 self-saturation mechanism (canonical at \kbleaf{research/2026-05-31\_Q-EMBED-SEL-1\_step\_c\_result.md}): the time-averaged phasor enclosed area at Axiom-4 self-saturation onset equals the Nyquist cell cross-section area, giving $\pi R r = \pi(d/2)^2 \Rightarrow R \cdot r = 1/4$ at $d = 1\,\ell_\text{node}$. Combined with regime (b) self-avoidance $R - r = 1/2$, this yields the Golden Torus geometry, at which $\Lambda_\text{surf} = 4\pi^2(R \cdot r) = \pi^2$ falls out automatically. The substrate-mechanism is universal across the $(2,q)$ ladder (electron $(2,3)$ + proton $(2,5)$ + $\Delta$ baryon $(2,7)$) and cross-domain (electron-scale self-saturation + nuclear-scale Pd-D externally-driven saturation).

Following the pattern of the PMNS angle derivation (Vol.~2, Ch.~3, \textit{Neutrino Sector Eigenvalues}), which rigorously establishes three mixing-angle values by identifying three distinct physical regimes with explicit boundary conditions, the $(2, 3)$ phase-space winding's geometric constraints on the Clifford torus partition into three regimes:

\paragraph{(a) Nyquist Regime (tube thickness).}
The discrete manifold has a fundamental sampling pitch $\ell_{node}$; modes of wavelength $< 2\ell_{node}$ are unresolvable (Axiom~1, Ch.~\ref{ch:fundamental_axioms}). The smallest topologically stable soliton is the $0_1$ unknot carrying the $(2, 3)$ phase-space winding (smallest non-trivial coprime $(p, q)$ pair), and the unknot's core tube must be resolvable, forcing $d = 1\,\ell_{node}$ as the rigid minimum diameter.

\paragraph{(b) Crossings Regime (transverse self-avoidance).}
At each topologically-marked crossing of the $(2, 3)$ phase-space winding, two strand portions of the bond-pair LC tank's $(V_{\text{inc}}, V_{\text{ref}})$ phasor trajectory pass in transverse proximity on the Clifford torus. The phase-space centerlines must be at least $d$ apart to avoid catastrophic dielectric rupture --- the same principle that sets the rigidity percolation threshold in the bulk (Ch.~\ref{ch:macroscopic_moduli}). For a Clifford-torus embedding this gives the closest centerline approach $2(R-r)$, hence $2(R-r) = d = 1 \Rightarrow R - r = 1/2$. This regime produces \emph{one} equation in $(R, r)$.

\paragraph{(c) Screening Regime (substrate Axiom-4 self-saturation at the (2,3) eigenmode).}
The substrate-mechanism for $R \cdot r = 1/4$ at $d = 1\,\ell_\text{node}$ closes via Axiom-4 self-saturation $+$ Op14 Meissner-asymmetric coupling at the $(2,3)$ chirality $+$ named phasor-area-equals-Nyquist-cell-area identification (canonical at \kbleaf{research/2026-05-31\_Q-EMBED-SEL-1\_step\_c\_result.md} \S2.3). The three substrate primitives that anchor the closure are:

\begin{enumerate}
\item \textbf{Bond LC tank reactive-energy capacity at saturation onset} (Ax 4 $+$ l3 \S3 Virial sum at \kbleaf{ave-kb/vol2/particle-physics/ch01-topological-matter/l3-electron-soliton-synthesis.md}:59): the maximum reactive energy a single K4 bond can store before Axiom-4 self-saturation engages is exactly $m_e c^2 = \hbar \omega_C$, distributed Virial-symmetrically between L and C sectors.
\item \textbf{Phasor enclosed area $=$ per-Compton-cycle reactive energy} (Op1 $+$ Op17 $+$ standard phasor analysis): the bond's per-cycle reactive energy is the closed-loop integral of $V\, dI$ around the phasor trajectory; for an elliptical trajectory with semi-axes $(R, r)$ in $(V_\text{inc}, V_\text{ref})$-coordinates in lattice-natural units, the enclosed area $\pi R r$ in $\ell_\text{node}^2$ units equals the bond's per-Compton-cycle reactive energy.
\item \textbf{Nyquist cell cross-section area} (Ax 1 $+$ regime (a) $+$ regime (b)): each K4 bond LC tank occupies one Nyquist cell with diameter $d = 1\,\ell_\text{node}$ and cross-section area $\pi(d/2)^2 = \pi/4$.
\end{enumerate}

The substrate-mechanism identification: at Axiom-4 self-saturation onset, the elliptical TIR envelope's phasor enclosed area equals the Nyquist cell's cross-section area (one quantum of bond reactive energy filling one Nyquist cell):
\begin{equation}
    \pi R r = \pi (d/2)^2 \quad \Rightarrow \quad \boxed{R \cdot r = (d/2)^2 = \tfrac{1}{4} \text{ at } d = 1\,\ell_\text{node}}
    \label{eq:screening_constraint}
\end{equation}

\paragraph{Note on the half-cover argument.} The standard-Clifford-torus half-cover argument for $\Lambda_{\text{surf}} = \pi^2$ --- casting it as a half-cover of the surface area $A_{\text{standard}} = 2\pi^2$ on $S^3$ at the canonical embedding $r_1 = r_2 = 1/\sqrt{2}$, with the spin-1/2 spinor-conjugate identification halving the physically-distinct observable surface --- is not load-bearing in the current derivation (it imports an SU(2) projective-ray identification). The substrate-native Axiom-4 self-saturation chain above carries the $\Lambda_{\text{surf}}$ value instead. The spin-1/2 structure of the electron substrate is derived via the K4 $\to$ A4 $\to$ 2T $\subset$ SU(2) $\to$ SO(3) chain $+$ Finkelstein-Misner mechanism (see \S\ref{sec:open_rigor}), independent of any surface-area-halving.

\paragraph{Substrate Origin of the $4\pi$ and $\pi^2$ Factors.}
The substrate $4\pi$ temporal-phase closure per observable Compton cycle (in $\Lambda_{\text{vol}}$) derives from the bipartite K4 lobe-count: per Compton period the soliton traverses both A-B sublattice lobes (factor 2) with phasor rotation $2\pi$ in each lobe (factor $2\pi$), giving $4\pi$ temporal-phase closure (canonical at \kbleaf{ave-kb/vol2/particle-physics/ch01-topological-matter/l3-electron-soliton-synthesis.md}:103-105: $m_e\text{(observable)} = m_\text{Cosserat}/2$ from bipartite K4 lobe traversal). The $\pi^2$ in $\Lambda_{\text{surf}}$ derives from the substrate-mechanism identification $R \cdot r = 1/4$ at Axiom-4 self-saturation onset (\kbleaf{research/2026-05-31\_Q-EMBED-SEL-1\_step\_c\_result.md} \S2.3). Both are substrate-native: the $4\pi$ via K4 bipartite lobe-count, the $\pi^2$ via Axiom-4 self-saturation phasor-area identification. The standard-physics-community names (``SU(2) $\to$ SO(3) double cover'', ``spinor half-cover'') are translation references to the substrate-mechanism content, not the substrate-mechanism content itself.

\paragraph{AVE-native provenance of the spin-1/2 structure (preserved as canonical for spin-1/2 structure).}
The spin-1/2 structure of the electron substrate is substrate-derived end-to-end via the K4 substrate's tetrahedral rotation symmetry:
\begin{enumerate}
    \item K4 rotation group: $T = A_4$ (proper tetrahedral rotation group, order $\lvert T \rvert = 12$; faithful representation on the 4-port basis verified per the K4-rotation-group derivation).
    \item Double cover: $2T \subset SU(2)$ (classical group theory; binary tetrahedral group as discrete SU(2) subgroup).
    \item Spin-$\tfrac{1}{2}$ emerges from the \emph{Finkelstein--Misner / Dirac-belt-trick mechanism} on the extended $0_1$ unknot defect embedded in the SO(3) manifold.
\end{enumerate}
Steps 1--3 fix the spin-1/2 structure of the electron substrate-natively. The half-cover surface-area identification is not load-bearing for $\Lambda_\text{surf} = \pi^2$; the $\pi^2$ value derives from the Axiom-4 self-saturation substrate-mechanism above (phasor-area-equals-Nyquist-cell-area at Axiom-4 self-saturation onset). The closure does not depend on ropelength-minimality embedding-selection on the canonical $r_1 = r_2 = 1/\sqrt{2}$ embedding (that framing is sidestepped, not invoked).\label{sec:open_rigor}

\paragraph{Closing the System.}
With three distinct physical regimes, we obtain three independent equations in the three unknowns $(R, r, d)$:
\begin{align}
    d &= 1 & &\text{(Nyquist, Axiom~1)} \nonumber \\
    2(R - r) &= d & &\text{(Crossings, Axiom~2 saturation at dielectric rupture)} \nonumber \\
    \pi R r &= \pi(d/2)^2 & &\text{(Screening, Axiom-4 self-saturation)} \nonumber
\end{align}
The system has the unique physical solution (Golden Torus, $R = \varphi/2$, $r = (\varphi-1)/2$) derived below. Each equation is grounded in a distinct physical principle, parallel to the three regime-boundary derivations that produce the three PMNS mixing angles. \emph{Honesty scoping (Class~B input-dependence):} the $R \cdot r = 1/4$ step (regime~(c) Axiom-4 self-saturation) routes through a substrate-canonical named identification (the phasor-area-equals-Nyquist-cell-area step per Q-EMBED-SEL-1), a canonical \emph{input} rather than a from-primitives derivation --- a Class~B substrate-mechanism manifestation, not a Class~2 emergence (canonical at \kbleaf{ave-kb/vol1/ch8-alpha-golden-torus.md} \S``Dual-axis classification'').

Solving this exact system of geometric hardware constraints ($R-r=1/2$ and $R\cdot r=1/4$) yields the exact physical bounding radii of the electron:
\begin{equation}
    R = \frac{1+\sqrt{5}}{4} = \frac{\varphi}{2} \approx 0.809 \quad \text{and} \quad r = \frac{-1+\sqrt{5}}{4} = \frac{\varphi-1}{2} \approx 0.309
    \label{eq:golden_torus_radii}
\end{equation}
Where $\varphi$ is the Golden Ratio. The electron is structurally locked not to an arbitrary heuristic, but to the \textbf{Golden Torus}---the absolute most mathematically compact non-intersecting geometry for a volume-bearing flux tube on a discrete grid.

\begin{figure}[htbp]
    \centering
    \includegraphics[width=0.9\textwidth]{trefoil_alpha_qfactor.png}
    \caption{\textbf{The Electron Soliton at Dielectric Ropelength.} The $0_1$ unknot's flux-tube geometry in real space, paired with the $(2,3)$ phase-space Clifford-torus winding, forces extreme flux crowding at the core, constrained by the discrete $\ell_{node}$ scale strictly to the Golden Torus limit ($R=\varphi/2$, $r=(\varphi-1)/2$) in phase space. Evaluating the Holomorphic Impedance at this absolute hardware boundary natively yields the geometric Q-factor ($\alpha^{-1} \approx 137.036$).}
    \label{fig:trefoil_soliton}
\end{figure}

\subsection{Derivation of $\Lambda_{\text{line}} = \pi \cdot d$ from Nyquist Quantization and Cross-Section 1-Cycle Integration}
\label{sec:lambda_line_rigor}

The line shape factor $\Lambda_{\text{line}} = \pi \cdot d$ rests on two ingredients drawn from the same regime structure that fixed $\Lambda_{\text{vol}}$ and $\Lambda_{\text{surf}}$: regime~(a) Nyquist quantization of the tube diameter, and the multipole-expansion identification of the 1-codimensional integral with the cross-section perimeter. The factor $\pi$ (not $2\pi$) is forced by the regime-(b) self-avoidance convention which makes $d$ a diameter rather than a radius. Each step is parallel in rigor (Class~B substrate-mechanism manifestation) to the substrate-mechanism derivation of regime~(c)'s $R \cdot r = 1/4$ in \S\ref{sec:pi_squared_rigor} (Axiom~4 self-saturation + Op14 Meissner-asymmetric coupling + named phasor-area-equals-Nyquist-cell-area identification per \kbleaf{research/2026-05-31\_Q-EMBED-SEL-1\_step\_c\_result.md}).

\paragraph{Identification of $\Lambda_{\text{line}}$ as the 1-cycle integral.}
The multipole decomposition on $\mathbb{T}^2 \subset S^3 \subset \mathbb{C}^2$ assigns one shape factor per geometric codimension: $\Lambda_{\text{vol}}$ to the 3-cycle phase volume, $\Lambda_{\text{surf}}$ to the 2-cycle Clifford-torus surface, and $\Lambda_{\text{line}}$ to the 1-cycle around the flux tube's transverse cross-section. On the LC network (Axiom~1), this 1-cycle integral is realized physically as the Amp\`ere loop $\oint \mathbf{B} \cdot d\boldsymbol{\ell}$ that links transverse flux to circulating displacement current --- the natural 1-D codimension integral on a Maxwell network. The line shape factor is the dimensionless length of this Amp\`ere loop.

\paragraph{Diameter convention from regime~(b).}
Regime~(b) self-avoidance imposes $2(R - r) = d$, where $2(R - r)$ is the closest centerline-to-centerline approach of the trefoil's strands at a topologically-marked crossing. For two flux tubes to just touch at their edges (without dielectric rupture), this centerline separation must equal the tube \emph{diameter}, not its radius. This makes $d$ unambiguously the tube diameter throughout the closure system; the corresponding tube radius is $d/2$.

\paragraph{Nyquist-quantized cross-section perimeter.}
Regime~(a) Nyquist quantizes the minimum lattice-resolvable tube diameter at $d = 1\,\ell_{\text{node}}$. Below this scale, the discrete LC network cannot represent any transverse cross-section structure at all --- sub-pitch ``shape'' distinctions are unphysical. The 1-cycle integral evaluated around the tube cross-section --- in the LC network's continuous-mode envelope, which is the only meaningful representation at the Nyquist resolving limit --- gives the standard perimeter form for a closed 1-cycle of diameter $d$:
\begin{equation}
    \Lambda_{\text{line}} \;=\; 2\pi \cdot (d/2) \;=\; \pi \cdot d
    \label{eq:lambda_line_form}
\end{equation}
This is the \emph{full} perimeter of a 1-cycle of diameter $d$ (equivalently $2\pi$ times its radius $d/2$). The factor $\pi$, rather than $2\pi$, is the necessary consequence of $d$ being expressed as a diameter (regime~(b) convention); there is no half-loop. The $\pi$ itself is the standard continuous angular factor of the closed 1-cycle integral, structurally parallel to the $2\pi$ factors that enter $\Lambda_{\text{vol}}$ and $\Lambda_{\text{surf}}$ through their 2- and 3-cycle integrals on $\mathbb{T}^2 \subset S^3 \subset \mathbb{C}^2$.

\paragraph{Closure.}
At the Nyquist-quantized minimum $d = 1\,\ell_{\text{node}}$:
\begin{equation}
    \Lambda_{\text{line}} \;=\; \pi \cdot d \;=\; \pi
    \label{eq:lambda_line_value}
\end{equation}
Each of the three multipole shape factors now carries the same structural form: a continuous angular factor (powers of $\pi$ from the closed cycles on $\mathbb{T}^2 \subset S^3 \subset \mathbb{C}^2$) multiplied by a regime-fixed scale. $R \cdot r = 1/4$ from regimes~(b)$\wedge$(c) fixes the scales of $\Lambda_{\text{vol}}$ and $\Lambda_{\text{surf}}$; $d = 1$ from regime~(a) Nyquist fixes the scale of $\Lambda_{\text{line}}$. The three regimes do uniform structural work across the three multipoles --- $\Lambda_{\text{line}}$ is no longer the odd one out. Numerical verification: \texttt{verify\_lambda\_line.py}.

\section{Multipole Decomposition of the Fine Structure Constant ($\alpha$) on $\mathbb{T}^2 \subset S^3 \subset \mathbb{C}^2$}
\label{sec:alpha_decomposition}

The Fine Structure Constant ($\alpha$) is not a randomly ``tuned'' magical scalar. It is identically the dimensionless topological self-impedance (Q-Factor) of this maximal-strain ground state. The total geometric impedance ($\alpha^{-1}$) is obtained by the Class~B mode-category decomposition of the Golden Torus's energy functional into its three codimensional Nyquist-cell categories (volumetric, surface, line) on $\mathbb{T}^2 \subset S^3 \subset \mathbb{C}^2$. The additive assembly is a \textbf{Class~B substrate-mechanism formalization} (Op21 multi-mode mode-counting at the $\Gamma = -1$ saturation boundary), not a Class~2 axiom-manifestation emergence: the Clifford-torus codimensional embedding remains a canonical input from upstream leaves.

\paragraph{Identification step.} The sum decomposition
\[
\alpha^{-1}_{\text{ideal}} \;\equiv\; \Lambda_{\text{vol}} \;+\; \Lambda_{\text{surf}} \;+\; \Lambda_{\text{line}}
\]
is the \emph{identification} of the three codimensional Nyquist-cell-count categories of the Clifford torus embedding (3-cell phase volume, 2-cell surface, 1-cell line) with the three independent substrate-mechanism contributions to the dimensionless self-impedance of the trefoil soliton. The cross-term-freeness of the addition is \textbf{Nyquist-cell / mode-category mutual exclusivity} at the $\Gamma = -1$ saturation boundary --- volume cells, surface cells, and line cells are mutually-exclusive Nyquist-cell categories and add without cross-terms --- \emph{not} parameter-orthogonality in $(R, r)$-space. Indeed the leaf's own forms give $\Lambda_{\text{vol}} = 16\pi^3(R \cdot r)$ and $\Lambda_{\text{surf}} = 4\pi^2(R \cdot r)$, so $(R \cdot r)$ is shared collinearly ($\Lambda_{\text{vol}}/\Lambda_{\text{surf}} = 4\pi$, the substrate temporal-phase closure per observable Compton cycle) and the $(R, r, d) \to \Lambda$ map has only 2-dimensional image; ``orthogonality'' therefore cannot mean $(R, r)$-independence. This is a \emph{multipole identification} at the level of mode-category mutual exclusivity, parallel in spirit to the multipole expansion of any extended electromagnetic source, and not an additional derivation step beyond the regime closure of \S\ref{sec:pi_squared_rigor}. Classification: \textbf{Class~B substrate-mechanism manifestation} (canonical walk-back at \kbleaf{ave-kb/vol1/ch8-alpha-golden-torus.md} \S``Schur-orthogonality status (Phase~3-A2 WALK-BACK closure, 2026-05-26)''; the Schur-orthogonality framing for the additive assembly is superseded by Op21 mode-counting at the saturation boundary).

Normalizing these limits by the fundamental spatial voxel ($\ell_{node}$) yields pure, dimensionless Impedance Shape Factors ($\Lambda$):

\begin{enumerate}
    \item \textbf{The Bulk (Volumetric Inductance, $\Lambda_{vol}$):} The hyper-volume of the 3-torus phase-space. The substrate $4\pi$ temporal-phase closure per observable Compton cycle derives from bipartite K4 lobe-count (2 sublattices) $\times\, 2\pi$ phasor rotation per lobe (canonical at \kbleaf{ave-kb/vol2/particle-physics/ch01-topological-matter/l3-electron-soliton-synthesis.md}:103-105). The standard-physics translation reference is ``SU(2) $\to$ SO(3) double-cover applied temporally'', but the substrate-mechanism content is K4 bipartite lobe-count, not an SU(2) postulate.
    \begin{equation}
        \Lambda_{vol} = (2\pi R)(2\pi r)(2\pi \cdot 2) = 16\pi^3 (R \cdot r) = 16\pi^3 \left(\frac{1}{4}\right) = \mathbf{4\pi^3} \approx 124.025
    \end{equation}

    \item \textbf{The Surface (Cross-Sectional Screening, $\Lambda_{surf}$):} The total geometric area of the Clifford Torus ($S^1 \times S^1$) bounding the knot.
    \begin{equation}
        \Lambda_{surf} = (2\pi R)(2\pi r) = 4\pi^2 (R \cdot r) = 4\pi^2 \left(\frac{1}{4}\right) = \mathbf{\pi^2} \approx 9.870
    \end{equation}

    \item \textbf{The Line (Tube Cross-Section Circumference, $\Lambda_{line}$):} The dimensionless circumference of the core flux-tube's cross-section evaluated at the minimum discrete node thickness ($d=1$). This is a 1D geometric line integral on the tube circumference, NOT a magnetic moment (which would have dimensions of area $\times$ current rather than length):
    \begin{equation}
        \Lambda_{line} = \pi \cdot d = \pi(1) = \mathbf{\pi} \approx 3.142
    \end{equation}
    Per the three-boundary-observables structure ($\mathcal{M}$, $\mathcal{Q}$, $\mathcal{J}$) canonical at the Common Foreword and KB leaf \texttt{ave-kb/common/boundary-observables-m-q-j.md}, this is the 1D-line projection of the $\alpha^{-1}$ three-integral decomposition, corresponding to the $\mathcal{Q}$ (boundary loop / linking number) dimensionality.
\end{enumerate}

Summing these three topological shape factors yields the closed-form theoretical invariant for a perfectly rigid ``Cold Vacuum'' (Absolute Zero, $0\,$K). Per the honest-$\alpha$ scope (Class B caveat below, \S Mathematical Closure): this is a \emph{closed-form geometric} invariant whose scale ($\sim 1/137$) is forced by the Compton-resonance trapping condition, but whose exact value rests on the one substrate-geometric identification $R \cdot r = 1/4$ that the substrate does not independently select --- Class B substrate-mechanism manifestation, not a parameter-free first-principles ``derivation'' of $\alpha$:

\begin{resultbox}{Cold-Lattice Fine-Structure Constant}
\begin{equation}
    \alpha^{-1}_{\text{ideal}} \equiv \Lambda_{vol} + \Lambda_{surf} + \Lambda_{line} = 4\pi^3 + \pi^2 + \pi \approx \mathbf{137.0363038}
    \label{eq:alpha_cold}
\end{equation}
\end{resultbox}

\textbf{Mathematical Closure (status):} The framework reduces 26 Standard Model parameters to a 3-element bounding set $\{m_e, \alpha, G\}$ + four axioms. The Layer~8 closure claim is that these three reduce to substrate geometry: (a) $\alpha^{-1}_{\text{ideal}} = 4\pi^3 + \pi^2 + \pi$ at the cold-lattice asymptote, from the Golden Torus three-regime closed form above --- whose scale ($\sim 1/137$) is forced by the Compton-resonance trapping condition but whose exact value rests on the one substrate-geometric identification $R \cdot r = 1/4$ that the substrate does not independently select (Class B substrate-mechanism manifestation; not a first-principles ``derivation,'' not zero-parameter); (b) the thermal running $\alpha^{-1}(T) = \alpha^{-1}_{\text{ideal}}(1 - \delta_{strain}(T))$ with $\delta_{strain}(T) > 0$ at $T > 0$ (existence and sign predicted; magnitude at $T_{\text{CMB}}$ currently fit to CODATA, see \S\ref{sec:alpha_thermal_running} predicted/fitted split); (c) one of $\{m_e, \ell_{node}\}$ computable from the other via the unknot ground state, with the other serving as the input mass scale. Axiom~1 (Ch.~\ref{ch:fundamental_axioms}) calibrates the baseline size of the lattice ($\ell_{node}$) to the electron rest-mass; this calibration is the input scale. Because the electron is the absolute structural failure mode of the lattice, its geometric packing Q-factor ($137.036$) \textbf{physically becomes} the macroscopic non-linear saturation limit for the rest of the universe (this is why $\alpha$ serves identically as the dielectric saturation bound in Axiom~4, \S\ref{sec:universal_saturation}).

Conditional on (a)--(c) at \textbf{Class B substrate-mechanism manifestation level}, the framework is zero-free-parameters at the SM level with $\alpha$ a \textbf{named geometric identification}: the $\delta_{strain}$ magnitude is a definitional residual (no derivation of its magnitude from substrate parameters exists in the corpus), and all three named $\alpha$-routes are closed-negative --- both golden-torus Class 2 lift-routes (dynamical selection $+$ the phasor$\leftrightarrow$real-space area bijection from K4 $+$ Cosserat primitives) stand at Class B, and the separate $z_0$-from-K4 (rigidity-percolation) route is closed-negative via the exact $z_{\text{eff}} \to 6$ Maxwell--Calladine count (\texttt{alpha\_free\_map\_to\_137\_exists = False}).
The cold-lattice asymptote derives from the Axiom-4 self-saturation substrate-mechanism (cross-particle and cross-domain universal) and the existence + sign of the thermal running are substrate-predicted; the magnitude of $\delta_{strain}$ at $T_{\text{CMB}}$ is one back-subtracted definitional residual (with the falsifiable claim that one scalar suffices across all temperatures). The practical count is ``26 SM parameters reduced to 3 bounding limits + 4 axioms; the $R \cdot r = 1/4$ relation in $\alpha$'s derivation closes at Class B substrate-mechanism manifestation via the Axiom-4 self-saturation substrate-mechanism (cross-particle and cross-domain universal); the named golden-torus Class 2 lift-routes (selection + bijection) stand at Class B; the thermal-running magnitude derivation of $\eta_\varepsilon$ within the Cosserat-rotation-sector mass-gap thermal-mode-population ASYM mechanism class (\kbleaf{vol3/cosmology/ch05-dark-sector/delta-strain-cosmic-tcc.md}) does not yield a magnitude --- the magnitude is a definitional residual, SIGN-only.''

Implementation: computed at engine import-time from pure geometric constants; see \texttt{ALPHA\_COLD\_INV} in \texttt{ave.core.constants}.

\textbf{Numerical verification (five scripts in \texttt{src/scripts/vol\_1\_foundations/}):}
\begin{itemize}
    \item \texttt{derive\_alpha\_from\_golden\_torus.py} --- evaluates the multipole decomposition at the Golden Torus, renders Fig.~\ref{fig:trefoil_soliton}, asserts agreement with \texttt{ALPHA\_COLD\_INV} to machine precision.
    \item \texttt{verify\_clifford\_half\_cover.py} --- historical 5-step verification of $\Lambda_{\text{surf}} = \pi^2$ via the spin-1/2 half-cover of the standard Clifford torus $\mathbb{T}^2 \subset S^3 \subset \mathbb{C}^2$. The script body preserves the algebraic identity at the Golden Torus geometry; the canonical substrate-mechanism derivation is at \kbleaf{research/2026-05-31\_Q-EMBED-SEL-1\_step\_c\_result.md} \S2.3.
    \item \texttt{verify\_lambda\_line.py} --- rigorously derives $\Lambda_{\text{line}} = \pi \cdot d$ from regime~(a) Nyquist quantization plus the regime-(b) diameter convention; verifies the multipole closure under varying $d$ and the structural symmetry of the three shape factors.
    \item \texttt{ropelength\_trefoil\_golden\_torus.py} --- minimizes the composite S$_{11}$-free-energy (ropelength + self-avoidance + Clifford-torus screening penalties) from arbitrary initial $(R, r)$; converges to $(\varphi/2, (\varphi-1)/2)$ to 7 decimal places as penalty weights $\to \infty$.
    \item \texttt{verify\_golden\_torus\_s11.py} --- ABCD cascade infrastructure + algebraic verification of the constraint system.
\end{itemize}
The derivation is algebraically closed and numerically grounded end-to-end.

\section{The Thermodynamic Expansion of Space (The Running Coupling)}
\label{sec:alpha_thermal_running}

The exact theoretical derivation yields $137.0363038$. However, the experimentally measured 2022 CODATA value is slightly lower: $\alpha^{-1}_{\text{exp}} \approx 137.035999$.

In the AVE framework, this discrepancy is not a mathematical error. It is a direct, measurable consequence of the \textbf{Thermal Expansion of the Universe}.

The ideal geometric value ($4\pi^3 + \pi^2 + \pi$) mathematically assumes a lattice with zero ambient kinetic energy. However, the physical universe is bathed in a thermodynamic heat bath: the Cosmic Microwave Background ($T \approx 2.7\,$K). Just as thermal energy physically expands a mechanical solid and lowers its elastic stiffness, the ambient heat of the universe physically expands the Cosserat vacuum, introducing stochastic phonon vibrations that fractionally soften its geometric impedance.

We natively define the Vacuum Strain Coefficient ($\delta_{strain}$) as this exact thermodynamic deviation:

\begin{resultbox}{Vacuum Strain Coefficient}
\begin{equation}
    \delta_{strain} = 1 - \frac{\alpha^{-1}_{\text{exp}}}{\alpha^{-1}_{\text{ideal}}} = 1 - \frac{137.035999}{137.036304} \approx \mathbf{2.22 \times 10^{-6}}
    \label{eq:delta_strain}
\end{equation}
\end{resultbox}

\noindent\textbf{Status disclosure (current edition; predicted/fitted split).}
\begin{itemize}
    \item \emph{Predicted (axiom-derived):} the cold-lattice asymptote $\alpha^{-1}_{\text{ideal}} = 4\pi^3 + \pi^2 + \pi \approx 137.0363038$; the existence of a positive thermal running of $\alpha^{-1}$ below this asymptote at any $T > 0$; the sign and the falsifiable claim that $\alpha^{-1}$ decreases further in regions of higher local thermal energy (collider cores, early universe).
    \item \emph{Fitted (one scalar at $T_{\text{CMB}}$):} the numerical magnitude $\delta_{strain} \approx 2.22\times 10^{-6}$, computed by back-subtraction from CODATA: $\delta_{strain} \equiv 1 - \alpha^{-1}_{\text{CODATA}}/\alpha^{-1}_{\text{ideal}}$, definitional given the engine's \texttt{DELTA\_STRAIN = 1 - (1/ALPHA)/ALPHA\_COLD\_INV} (\texttt{src/ave/core/constants.py}). The attribution to spatial-metric thermal expansion at $T_{\text{CMB}} = 2.725\,$K is a physical narrative consistent with the predicted sign; it is not yet a derivation of the magnitude from $T_{\text{CMB}}$ + a lattice elastic modulus.
    \item \emph{Tested:} that \textbf{one} thermal scalar suffices to bridge the cold-lattice asymptote to the CODATA value. Multiple-temperature measurements of $\alpha$ (collider cores, primordial nucleosynthesis $\alpha$, ultracold cavity experiments) test the same one-scalar structural claim at different $T$.
\end{itemize}
This is the same disclosure pattern as Vol~6 (one fitted scalar per nucleus, with predicted topology + parameter count): structure predicted, magnitude fitted, falsifiable across the relevant variable. To upgrade $\delta_{strain}$ from a one-scalar fit to a true zero-parameter prediction, the chain needs to compute the lattice's effective thermal expansion coefficient from first principles --- the bulk modulus $G_{vac}$ (Ch.~\ref{ch:electrodynamics}) and the equipartition energy at $T_{\text{CMB}}$ --- and verify that the predicted $\delta_{strain}$ matches the back-subtracted value within tolerance. That magnitude-derivation does not currently appear in the corpus.

Implementation: see \texttt{DELTA\_STRAIN} in \texttt{ave.core.constants}, computed from CODATA \texttt{ALPHA} divided by \texttt{ALPHA\_COLD\_INV}.

\textbf{Falsifiable Prediction:} Because $\alpha$ is defined as a literal mechanical property of a physical lattice, it must act as a \textit{Running Coupling Constant}. If measured in a region of extreme localized thermal energy (e.g., inside a particle collider, or in the hot early universe), the localized stress will dynamically expand the lattice bonds, causing $\alpha^{-1}$ to decrease further. Conversely, the ideal theoretical limit $137.0363038$ serves as the exact impenetrable mathematical asymptote at true absolute zero.

This is the third distinct thermal running of $\alpha$ predicted by AVE, distinct from:
\begin{itemize}
    \item \textbf{Standard QED vacuum polarization} (which $\alpha$ is conventionally modeled to exhibit above the electron-positron threshold): different mechanism, different scale.
    \item \textbf{Proton thermal softening} ($\delta_{th} = 1/(14\pi^2) \approx 7.21\times 10^{-3}$): the thermal correction at the proton's $10^{13}$\,K core that enters the Faddeev-Skyrme coupling $\kappa_{FS}$. A different temperature regime and a different physical quantity (the Skyrme coupling, not $\alpha$ directly).
\end{itemize}

The Vacuum Strain Coefficient is the $\alpha$-specific, CMB-sourced running that bridges the pure geometric cold prediction to the observed value.
